\documentclass[12pt]{article}

\usepackage[utf8]{inputenc}
\usepackage{amsmath, amssymb, amsthm}
\usepackage{geometry}
\geometry{a4paper, margin=2.5cm}

\title{Teorema Aplicației Deschise}
\author{}
\date{}

\theoremstyle{definition}
\newtheorem{definition}{Definiție}

\theoremstyle{plain}
\newtheorem{theorem}{Teoremă}
\newtheorem{proposition}{Propoziție}

\begin{document}

\maketitle

\begin{definition}
Fie $X$ și $Y$ două spații Banach. Un operator liniar continuu


\[
T : X \to Y
\]


se numește \emph{surjectiv} dacă pentru orice $y \in Y$ există $x \in X$ cu $T(x) = y$.
\end{definition}

\begin{theorem}[Teorema Aplicației Deschise]
Fie $X$ și $Y$ două spații Banach și $T : X \to Y$ un operator liniar continuu și surjectiv. Atunci $T$ este o aplicație deschisă, adică pentru orice mulțime deschisă $U \subset X$, imaginea $T(U)$ este deschisă în $Y$.

Echivalent, există o constantă $c > 0$ astfel încât


\[
T(B_X(0,1)) \supset B_Y(0,c),
\]


unde $B_X(0,1)$ și $B_Y(0,c)$ sunt bilele deschise de rază $1$, respectiv $c$.
\end{theorem}

\begin{proposition}
În ipotezele teoremei aplicației deschise, următoarele afirmații sunt echivalente:
\begin{enumerate}
    \item[(i)] $T$ este surjectiv.
    \item[(ii)] Există $c > 0$ astfel încât pentru orice $y \in Y$ există $x \in X$ cu $T(x) = y$ și
    

\[
    \|x\| \le c \|y\|.
    \]


    \item[(iii)] Imaginea lui $T$ conține o bilă deschisă nenulă.
\end{enumerate}
\end{proposition}

\end{document}
